\section{AI-Transformed Soundscapes as a Design Instance}

The final part of the dissertation develops one design instance informed by the architecture and spatial-expert studies. It focuses on sound as a medium through which affective and biophilic interaction can be explored.

Sound is central to how shared indoor environments are experienced. Speech, typing, footsteps, ventilation, equipment, doors, and sudden impacts provide information about activity and the presence of others. These sounds may support social awareness, but they may also interrupt concentration, reduce privacy, or contribute to fatigue. Indoor sound is therefore commonly treated as a problem of noise reduction, acoustic separation, or masking.

Conventional nature-sound interventions generally add recordings of birds, water, wind, or other natural environments to a room. Such sounds may create a more pleasant atmosphere or reduce the prominence of unwanted noise. They nevertheless remain separate from the activity of the space. The workplace produces one soundscape, while the nature recording provides another layer over it.

The prototype developed in this dissertation examines a different proposition. It uses artificial intelligence to transform ongoing office sound into nature-like auditory textures. Rather than placing a fixed nature recording over the environment, the system treats office sound as the input for a responsive biophilic soundscape. Changes in office activity affect the transformed sound, while the original sources become less directly recognisable.

The prototype examines whether everyday office sound can be reframed rather than simply removed or masked. It also considers whether environmental responsiveness can be communicated through ambience rather than through an explicit display or notification. In this system, artificial intelligence is used as a transformational medium rather than as a means of classifying occupants, inferring emotional states, or selecting an optimal response.

The aim is not to propose a universal solution to office noise. The system functions as a design instance through which affective, multisensory, social, and biophilic possibilities can be examined. It makes the broader thesis argument available in a concrete form that participants and experts can experience, interpret, and question.

The prototype is studied through HCI methods, including an online participant study and expert discussions. These studies examine how the transformed sound is described, whether it is perceived as nature-like, how it differs from conventional nature recordings, whether it appears connected to office activity, and what forms of value or concern it produces. They also address control, distraction, artificiality, privacy, spatial integration, and sustained use.

Through this final study, the dissertation returns from design possibility to empirical investigation. The prototype demonstrates one way in which affective understanding and biophilic design can inform an interactive indoor environment, while its evaluation identifies the limitations and tensions that future work must address.